% 英文摘要：标题 Arial 小三居中；正文 Times New Roman 小四，固定行距 20pt
% 单独成页，不出现图表
\xuptheadingnotoc{\sffamily\fontsize{15pt}{20pt}\selectfont ABSTRACT}

With the continuous evolution of network attacks, traditional rule-based anomaly detection methods can hardly balance the detection rate and the false alarm rate. Aiming at the problems of high-dimensional traffic data and scarce attack samples in a campus network environment, this paper designs and implements a deep learning based network anomaly detection system. The system constructs input vectors from traffic statistical features, learns the latent distribution of normal traffic with a deep autoencoder, and identifies anomalous behaviors by reconstruction errors. An attention mechanism is further introduced to improve the sensitivity of the model to key features. Experimental results show that the proposed method achieves higher accuracy and recall than competing methods on public datasets, and meets the requirement of real-time detection in practical scenarios.

In terms of system design, this paper completes the design and implementation of the data acquisition module, feature engineering module, model training module and visual alarm module. The modules are loosely coupled for further extension. Finally, functional and performance tests are carried out, and the results show that the system runs stably and the detection delay meets the design requirements.

\vspace{0.5em}
\noindent{\bfseries\zihao{-4} Key words：}deep learning；network anomaly detection；autoencoder；traffic analysis；feature engineering
